% ===================================================================
%  TAULA N.º 16 — Els Grans Magatzems. — El Basar. — Les Joguines.
%  Traduction 2026 d'apres texto/io, controlee sur texto/fr et
%  texto/es. Consigne : voir texto/ca/10-taula-01.tex.
%
%  Le tableau d'astronomie (bloc 15-1) porte des renvois a LETTRES,
%  (a) a (f), graves sur le tableau lui-meme. L'ido saute la lettre
%  (d), que le fac-simile grave pourtant : la colonne la rend, comme
%  toutes les autres traductions. Ecart declare dans outils/renvoji.py.
%
%  DERNIER TABLEAU DE LA COLONNE CATALANE. Bilan des seize : zero
%  inversion. C'est la troisieme colonne dans ce cas, apres
%  l'esperanto et l'interlingua, et la premiere qui soit une langue
%  naturelle. La raison est la meme pour les trois : l'ido a pris son
%  ordre des mots aux langues romanes, et le catalan est l'une
%  d'elles.
% ===================================================================

%%K t16-apar-1 apar
\VUsaut{4.0mm}
\VUtitre{15.0pt}{60}{TAULA N.º 16}
\VUsaut{2.0mm}
\VUpk{13.2pt}{40}{Els Grans Magatzems. --- El Basar.}
\VUpk{13.2pt}{40}{Les Joguines.}
\VUsaut{3.4mm}

%%K t16-01-1 p
1. --- S'acosta l'any nou. Els grans magatzems han preparat les seves exposicions
de \VUgras{joguines} \textsuperscript{(1)} i regals d'any nou.

%%K t16-01-2 p
Vaig demanar al meu avi que em portés al basar a veure aquella bella exposició, i
va acceptar.

%%K t16-02-1 p
2. --- Primer ens vam aturar davant d'unes belles panòplies d'armes amb un
\VUgras{fusell}, una \VUgras{gorra de plat}, un \VUgras{sabre}, un \VUgras{tahalí}
i un parell de \VUgras{xarreteres}, o bé un \VUgras{casc}, una \VUgras{cuirassa}
\textsuperscript{(3)} i una \VUgras{pistola} \textsuperscript{(4)}. Tot allò
m'interessava molt més que les \VUgras{llanternes màgiques}
\textsuperscript{(2)}.

%%K t16-tit-1 sub
\VUsaut{3.0mm}
\VUpk{10.2pt}{40}{Bells espectacles.}
\VUsaut{2.6mm}

%%K t16-03-1 p
3. --- A la mateixa secció hi havia un \VUgras{sentinella} \textsuperscript{(5)} a
la seva \VUgras{garita}; semblava muntar guàrdia damunt tot un parc de feres de
cartró. Quants animals hi havia! Un \VUgras{lloro} \textsuperscript{(6)} al seu
bastó, un \VUgras{rinoceront} \textsuperscript{(7)} amb una banya al nas, un
\VUgras{ren} \textsuperscript{(8)}, un \VUgras{estruç} \textsuperscript{(9)}, un
\VUgras{mico} \textsuperscript{(10)} que trencava una nou, una \VUgras{guineu}
\textsuperscript{(11)} que s'enduia una gallina, un \VUgras{cocodril}
\textsuperscript{(12)} amb les seves enormes barres obertes, un \VUgras{camell}
\textsuperscript{(13)}, un elegant \VUgras{llebrer} \textsuperscript{(14)}, una
\VUgras{pantera} \textsuperscript{(15)}, i després dos bitxos que jo no coneixia,
que són, segons em diuen, una \VUgras{mostela} \textsuperscript{(16)} i una
\VUgras{hiena} \textsuperscript{(17)}. Hi havia també un \VUgras{lleopard}
\textsuperscript{(18)} i un \VUgras{llop} \textsuperscript{(21)}; molt a prop hi
havia bufones \VUgras{rates} \textsuperscript{(19)} i \VUgras{ratolins}
\textsuperscript{(20)} diminuts de cautxú.

%%K t16-03-1 p suite
Finalment, un \VUgras{hipopòtam} \textsuperscript{(22)} i un \VUgras{dromedari}
\textsuperscript{(23)} geperut completaven la col·lecció. M'hi hauria quedat hores
enteres si no hi hagués hagut, molt a prop, un esplèndid campament ple de
\VUgras{soldadets de plom} \textsuperscript{(29)} que em va cridar l'atenció.

%%K t16-tit-2 sub
\VUsaut{4.0mm}
\VUpk{10.2pt}{40}{Soldats i guerra.}
\VUsaut{2.8mm}

%%K t16-04-1 p
4. --- El meu avi, que és un \VUgras{invàlid de guerra} \textsuperscript{(24)} i va
haver de deixar l'exèrcit arran d'una \VUgras{ferida} que li va valer la
\VUgras{medalla militar}, adora explicar-me les \VUgras{campanyes} en què va
participar. Li va encantar explicar-me els diferents moviments del petit exèrcit
en miniatura. Un grup de soldats de debò que visitaven el basar --- un
\VUgras{enginyer militar} \textsuperscript{(25)}, un \VUgras{caçador alpí}
\textsuperscript{(26)} amb la seva \VUgras{boina}, un \VUgras{sergent primer}
\textsuperscript{(27)} d'infanteria i un \VUgras{sergent} \textsuperscript{(28)}
d'artilleria amb un \VUgras{galó} d'or a la mànega --- es va aturar vora nostre a
escoltar les explicacions.

%%K t16-05-1 p
5. --- El meu avi, molt orgullós de tenir un públic tan distingit, va improvisar
tot un pla de batalla.

%%K t16-05-2 p
La plaça forta, amb les seves \VUgras{muralles} i els seus \VUgras{fossats},
estava assetjada per l'\VUgras{exèrcit enemic}, que, després d'un llarg
\VUgras{bombardeig}, es disposava a prendre-la per \VUgras{assalt}. Però els
\VUgras{assetjats}, empesos per la fam, havien intentat una \VUgras{sortida}
contra el \VUgras{campament} dels \VUgras{assetjadors}. Rebutjat l'\VUgras{atac}
en un aferrissat \VUgras{combat} al voltant de les \VUgras{tendes}, els
assetjadors \VUgras{vencedors} van llançar els seus \VUgras{esquadrons} en
\VUgras{persecució} dels assaltants \VUgras{vençuts}. Aquests \VUgras{es van
retirar}, deixant el \VUgras{camp de batalla} sembrat de \VUgras{morts} i
\VUgras{ferits}. Els \VUgras{llitericers} van portar els ferits a l'\VUgras{hospital
de campanya} perquè els curés el \VUgras{cirurgià}.

%%K t16-05-3 p
Malgrat l'heroica resistència d'uns quants \VUgras{batallons} que s'havien format
en \VUgras{quadre}, la \VUgras{derrota} va ser completa; la resta de l'exèrcit va
\VUgras{fugir}, deixant enrere molts \VUgras{presoners}. L'enemic estava a punt de
coronar la seva \VUgras{victòria} prenent la plaça.

%%K t16-05-3 p suite
Potser fóra aquell el final de la campanya, i després d'un \VUgras{armistici}
podria signar-se un \VUgras{tractat de pau}.

%%K t16-tit-3 sub
\VUsaut{4.4mm}
\VUpk{10.2pt}{40}{Els regals d'any nou.}
\VUsaut{2.8mm}

%%K t16-06-1 p
6. --- Els joves soldats, que reien escoltant el pintoresc relat del veterà, li
van fer algunes preguntes més. Quant a mi, vaig fer la volta al taulell per mirar
una bella \VUgras{ballesta} \textsuperscript{(32)} i un \VUgras{arc}
\textsuperscript{(33)} amb la seva \VUgras{fletxa} \textsuperscript{(31)}. Allà
vaig trobar una altra col·lecció d'animals: dos \VUgras{ocells cantors}
\textsuperscript{(34)}, un \VUgras{rossinyol} i una \VUgras{tallareta} de corda que
cantaven a tot pulmó, i una sèrie de bitxos estranys --- un \VUgras{ratpenat}
\textsuperscript{(35)}, una \VUgras{boa} \textsuperscript{(36)}, una
\VUgras{tortuga} \textsuperscript{(37)}, una \VUgras{foca} \textsuperscript{(38)},
una \VUgras{balena} \textsuperscript{(39)} i un \VUgras{tauró}
\textsuperscript{(40)} fets de pell daurada.

%%K t16-07-1 p
7. --- No em vaig aturar gaire a mirar-los, car havia vist el que somiava: un
magnífic \VUgras{teatre de titelles} \textsuperscript{(41)} amb esplèndids
\VUgras{decorats} i un deliciós \VUgras{teló} vermell. A l'\VUgras{escenari}, dos
actors semblaven representar un \VUgras{paper} en una \VUgras{obra}
interessantíssima, car els petits \VUgras{espectadors} de cartró de les llotges,
o asseguts a les \VUgras{butaques} i al \VUgras{galliner} darrere del
\VUgras{director d'orquestra}, aplaudien amb entusiasme.

%%K t16-07-2 p
Per desgràcia, aquell teatre era caríssim.

%%K t16-08-1 p
8. --- A més, era difícil de triar entre les joguines, car els aparadors eren
plens de joguines de tota mena: \VUgras{Arlequí} \textsuperscript{(42)},
\VUgras{Polxinel·la} \textsuperscript{(43)}, \VUgras{Pierrot}
\textsuperscript{(44)}, \VUgras{òlibes} \textsuperscript{(45)} que batien les ales,
\VUgras{merles} \textsuperscript{(46)} que xiulaven, \VUgras{gossos d'aigües} que
lladraven, un \VUgras{gramòfon} \textsuperscript{(47)} i \VUgras{trencaclosques}.

%%K t16-noto-1 noto
\VUnotes{143.2mm}{%
\textit{(*) Vegeu la taula n.º 6.}%
}

%%K t16-08-1 p suite
Una d'aquelles joguines em va divertir moltíssim: representava un \VUgras{combat
naval} \textsuperscript{(48)} en què un \VUgras{vaixell} és atacat per
\VUgras{pirates} davant d'una costa plantada de \VUgras{cocoters}.

%%K t16-09-1 p
9. --- Vaig poder mirar totes aquestes meravelles amb tota calma, car hi havia
poca gent als magatzems --- amb prou feines un grapat de curiosos, un
\VUgras{dragó} \textsuperscript{(49)} i un \VUgras{cobrador}
\textsuperscript{(50)} que presentava una \VUgras{lletra de canvi} a la
\VUgras{caixa} \textsuperscript{(51)} principal.

%%K t16-10-1 p
10. --- Vaig preguntar al meu avi per a què servien els llibres que hi havia als
prestatges darrere del \VUgras{caixer}; em va dir que eren els \VUgras{llibres de
comptes}.

%%K t16-tit-4 sub
\VUsaut{3.6mm}
\VUpk{10.2pt}{40}{Mirant per tota la botiga.}
\VUsaut{2.8mm}

%%K t16-11-1 p
11. --- En sortir de la secció de joguines, vam anar a la basteria a fer un cop
d'ull a les \VUgras{selles de muntar} \textsuperscript{(53)}, els \VUgras{estreps}
\textsuperscript{(54)}, les \VUgras{brides} \textsuperscript{(55)}, els
\VUgras{frens} \textsuperscript{(56)} i les \VUgras{fuetes}. Mentre el \VUgras{cap
de secció} \textsuperscript{(57)} donava al meu avi algunes
informacions, vaig anar a mirar l'exposició de \VUgras{porcellana} i
\VUgras{cristalleria} \textsuperscript{(58)}, on hi havia belles
\VUgras{enciameres} i delicades \VUgras{teteres} \textsuperscript{(59)}.

%%K t16-12-1 p
12. --- Per pujar al primer pis, vam passar per darrere de l'aparador de les
\VUgras{cotilles} \textsuperscript{(60)}, vora la merceria, on hi havia exposats
bellíssims \VUgras{calçotets} \textsuperscript{(63)}. A prop, una
\VUgras{dependenta} \textsuperscript{(61)} ajudava una \VUgras{clienta}
\textsuperscript{(62)} a escollir belles \VUgras{teles} \textsuperscript{(64)} de
seda.

%%K t16-12-2 p
En pujar l'escala, vaig notar que els magatzems estaven adornats amb
\VUgras{fanalets} \textsuperscript{(65)} japonesos, \VUgras{ombrel·les}
\textsuperscript{(66)} i enormes \VUgras{palmeres} \textsuperscript{(70)}.

%%K t16-13-1 p
13. --- Al primer pis no hi havia gaire cosa a veure: només mobles (*),
\VUgras{caixes de cabals} \textsuperscript{(67)}, \VUgras{calaixeres}
\textsuperscript{(68)} i \VUgras{cotxets de nen} \textsuperscript{(69)}. Però la
vista general dels magatzems era molt \VUgras{entretinguda}.

%%K t16-14-1 p
14. --- Sota nostre veia els instruments de música: \VUgras{arpes}
\textsuperscript{(71)}, \VUgras{guitarres} \textsuperscript{(72)},
\VUgras{mandolines} \textsuperscript{(73)}, \VUgras{cornetins}
\textsuperscript{(74)}, \VUgras{clarinets} \textsuperscript{(75)}, violins amb els
seus \VUgras{arcs} \textsuperscript{(76)} i fins i tot un \VUgras{bombo}
\textsuperscript{(77)}. Una mica més enllà, al taulell de papereria, un
\VUgras{estudiant} \textsuperscript{(78)} regatejava amb el \VUgras{dependent}
\textsuperscript{(79)} una cartera de documents. Vora d'ells distingia un bell
\VUgras{abecedari} \textsuperscript{(80)}.

%%K t16-14-2 p
Al mig dels magatzems s'alçava un alt \VUgras{arbre de Nadal}
\textsuperscript{(81)} carregat d'espelmes, galindaines i joguines de tota mena:
\VUgras{cavalls de fusta} \textsuperscript{(82)}, \VUgras{nines}
\textsuperscript{(83)}, etc.

%%K t16-14-3 p
La \VUgras{perfumeria} \textsuperscript{(84)}, amb els seus \VUgras{dentifricis} i
els seus diversos \VUgras{perfums}, escampava una olor molt agradable que pujava
fins a nosaltres.

%%K t16-tit-5 sub
\VUsaut{3.4mm}
\VUpk{10.2pt}{40}{Comprem alguna cosa.}
\VUsaut{2.8mm}

%%K t16-15-1 p
15. --- Com que el meu avi començava a trobar el temps llarg, vam tornar a baixar
per la gran escala. Es va aturar un moment davant d'un \VUgras{quadre
d'astronomia} \textsuperscript{(85)} que representava, segons em va dir,
\VUgras{cometes} \textsuperscript{(a)}, \VUgras{eclipsis de lluna i de sol}
\textsuperscript{(b)}, una rosa dels vents amb els \VUgras{quatre punts cardinals}
\textsuperscript{(c)} \VUgras{(nord, sud, est i oest)}, un mapa dels dos
\VUgras{hemisferis} \textsuperscript{(d)}, els \VUgras{planetes}
\textsuperscript{(e)} i el \VUgras{sistema solar} \textsuperscript{(f)}. No vaig
entendre res d'allò, i em van interessar molt més la lliurea galonejada d'un
\VUgras{repartidor} \textsuperscript{(86)} que passava, i els bonics
\VUgras{ventalls} \textsuperscript{(87)} que tenia al costat, vora els
\VUgras{portamonedes} \textsuperscript{(88)} i les \VUgras{carteres}
\textsuperscript{(89)}.

%%K t16-16-1 p
16. --- Vaig admirar també al taulell unes \VUgras{plomes}
\textsuperscript{(90)} i unes \VUgras{sivelles de cinturó} \textsuperscript{(91)},
i les boniques \VUgras{flors artificials} \textsuperscript{(94)} graciosament
disposades en cistelles.

%%K t16-16-1 p suite
Les \VUgras{violes}, els \VUgras{alelís}, els \VUgras{jacints}
\textsuperscript{(95)}, els \VUgras{geranis} \textsuperscript{(96)}, els
\VUgras{lliris} \textsuperscript{(97)} i les \VUgras{caputxines}
\textsuperscript{(98)} estaven molt ben imitats.

%%K t16-tit-6 sub
\VUsaut{4.4mm}
\VUpk{10.2pt}{40}{El regal de l'avi.}
\VUsaut{2.8mm}

%%K t16-17-1 p
17. --- Vaig demanar al meu avi que em comprés un dels \VUgras{raspalls de dents}
\textsuperscript{(92)} d'una caixa que hi havia vora les \VUgras{agulles de cap}
\textsuperscript{(93)}. Va acceptar, i vaig anar jo mateix a pagar la meva compra
a la caixa, al costat de la qual estaven fent un gran \VUgras{paquet}
\textsuperscript{(100)} per a un \VUgras{client} \textsuperscript{(99)}.

%%K t16-18-1 p
18. --- De sobte hi va haver un lleuger enrenou entre la gent que era vora els
\VUgras{bronzes} \textsuperscript{(101)} artístics. Un \VUgras{lladre}
\textsuperscript{(102)} acabava de ser sorprès en flagrant delicte per un
\VUgras{inspector dels magatzems} \textsuperscript{(103)} quan intentava robar
algú. Es va manar buscar un guàrdia per \VUgras{detenir-lo}, i en el moment en
què sortíem, s'enduien el culpable a la comissaria.
\VUsaut{6.0mm}
\VUfilet{22mm}
